
\documentclass[a4paper,11pt]{article}
\usepackage[a4paper, margin=8em]{geometry}

% usa i pacchetti per la scrittura in italiano
\usepackage[french,italian]{babel}
\usepackage[T1]{fontenc}
\usepackage[utf8]{inputenc}
\frenchspacing 

% usa i pacchetti per la formattazione matematica
\usepackage{amsmath, amssymb, amsthm, amsfonts}

% usa altri pacchetti
\usepackage{gensymb}
\usepackage{hyperref}
\usepackage{standalone}

% imposta il titolo
\title{Appunti Elettronica Digitale}
\author{Luca Seggiani}
\date{2026}

% imposta lo stile
% usa helvetica
\usepackage[scaled]{helvet}
% usa palatino
\usepackage{palatino}
% usa un font monospazio guardabile
\usepackage{lmodern}

\renewcommand{\rmdefault}{ppl}
\renewcommand{\sfdefault}{phv}
\renewcommand{\ttdefault}{lmtt}

% disponi il titolo
\makeatletter
\renewcommand{\maketitle} {
	\begin{center} 
		\begin{minipage}[t]{.8\textwidth}
			\textsf{\huge\bfseries \@title} 
		\end{minipage}%
		\begin{minipage}[t]{.2\textwidth}
			\raggedleft \vspace{-1.65em}
			\textsf{\small \@author} \vfill
			\textsf{\small \@date}
		\end{minipage}
		\par
	\end{center}

	\thispagestyle{empty}
	\pagestyle{fancy}
}
\makeatother

% disponi teoremi
\usepackage{tcolorbox}
\newtcolorbox[auto counter, number within=section]{theorem}[2][]{%
	colback=blue!10, 
	colframe=blue!40!black, 
	sharp corners=northwest,
	fonttitle=\sffamily\bfseries, 
	title=~\thetcbcounter: #2, 
	#1
}

% disponi definizioni
\newtcolorbox[auto counter, number within=section]{definition}[2][]{%
	colback=red!10,
	colframe=red!40!black,
	sharp corners=northwest,
	fonttitle=\sffamily\bfseries,
	title=~\thetcbcounter: #2,
	#1
}

% U.D.M
\newcommand{\amp}{\ensuremath{\mathrm{A}}}
\newcommand{\volt}{\ensuremath{\mathrm{V}}}
\newcommand{\meter}{\ensuremath{\mathrm{m}}}
\newcommand{\second}{\ensuremath{\mathrm{s}}}
\newcommand{\farad}{\ensuremath{\mathrm{F}}}
\newcommand{\henry}{\ensuremath{\mathrm{H}}}
\newcommand{\siemens}{\ensuremath{\mathrm{S}}}

% circuiti
\usepackage{circuitikz}
\usetikzlibrary{babel}

% disegni
\usepackage{pgfplots}
\pgfplotsset{width=10cm,compat=1.9}

% disponi codice
\usepackage{listings}
\usepackage[table]{xcolor}

\definecolor{codegreen}{rgb}{0,0.6,0}
\definecolor{codegray}{rgb}{0.5,0.5,0.5}
\definecolor{codepurple}{rgb}{0.58,0,0.82}
\definecolor{backcolour}{rgb}{0.95,0.95,0.92}

\lstdefinestyle{codestyle}{
		backgroundcolor=\color{black!5}, 
		commentstyle=\color{codegreen},
		keywordstyle=\bfseries\color{magenta},
		numberstyle=\sffamily\tiny\color{black!60},
		stringstyle=\color{green!50!black},
		basicstyle=\ttfamily\footnotesize,
		breakatwhitespace=false,         
		breaklines=true,                 
		captionpos=b,                    
		keepspaces=true,                 
		numbers=left,                    
		numbersep=5pt,                  
		showspaces=false,                
		showstringspaces=false,
		showtabs=false,                  
		tabsize=2
}

\lstdefinestyle{shellstyle}{
		backgroundcolor=\color{black!5}, 
		basicstyle=\ttfamily\footnotesize\color{black}, 
		commentstyle=\color{black}, 
		keywordstyle=\color{black},
		numberstyle=\color{black!5},
		stringstyle=\color{black}, 
		showspaces=false,
		showstringspaces=false, 
		showtabs=false, 
		tabsize=2, 
		numbers=none, 
		breaklines=true
}

\lstdefinelanguage{javascript}{
	keywords={typeof, new, true, false, catch, function, return, null, catch, switch, var, if, in, while, do, else, case, break},
	keywordstyle=\color{blue}\bfseries,
	ndkeywords={class, export, boolean, throw, implements, import, this},
	ndkeywordstyle=\color{darkgray}\bfseries,
	identifierstyle=\color{black},
	sensitive=false,
	comment=[l]{//},
	morecomment=[s]{/*}{*/},
	commentstyle=\color{purple}\ttfamily,
	stringstyle=\color{red}\ttfamily,
	morestring=[b]',
	morestring=[b]"
}

% disponi sezioni
\usepackage{titlesec}

\titleformat{\section}
	{\sffamily\Large\bfseries} 
	{\thesection}{1em}{} 
\titleformat{\subsection}
	{\sffamily\large\bfseries}   
	{\thesubsection}{1em}{} 
\titleformat{\subsubsection}
	{\sffamily\normalsize\bfseries} 
	{\thesubsubsection}{1em}{}

% disponi alberi
\usepackage{forest}

\forestset{
	rectstyle/.style={
		for tree={rectangle,draw,font=\large\sffamily}
	},
	roundstyle/.style={
		for tree={circle,draw,font=\large}
	}
}

% disponi algoritmi
\usepackage{algorithm}
\usepackage{algorithmic}
\makeatletter
\renewcommand{\ALG@name}{Algoritmo}
\makeatother

% disponi numeri di pagina
\usepackage{fancyhdr}
\fancyhf{} 
\fancyfoot[L]{\sffamily{\thepage}}

\makeatletter
\fancyhead[L]{\raisebox{1ex}[0pt][0pt]{\sffamily{\@title \ \@date}}} 
\fancyhead[R]{\raisebox{1ex}[0pt][0pt]{\sffamily{\@author}}}
\makeatother

\begin{document}
% sezione (data)
\section{Lezione del 15-05-26}

% stili pagina
\thispagestyle{empty}
\pagestyle{fancy}

% testo

\subsection{Multivibratori}
I \textbf{multivibratori} in elettronica sono una serie di circuiti, divisa in:
\begin{itemize}
	\item Multivibratori \textbf{monostabili}: rimangono in uno stato stabile se non sollecitati, quindi se sollecitati (da un \textit{trigger}) entrano in uno secondo stato dove rimangono per un tempo $T$, e da cui successivamente decadono allo stato stabile. L'esempio classico è quello del timer; 
	\item Multivibratori \textbf{bistabili}: hanno 2 stati stabili, e il passaggio da uno stato all'altro avviene in presenza di una data sollecitazione. Formano i \textit{latch} e i \textit{flip-flop}; 
	\item Multivibratori \textbf{astabili}: questi sono in uno stato continuamente instabile, e vengono quindi detti anche \textit{oscillatori}. Un esempio è il generatore del segnale di \textit{clock}.
\end{itemize}

\subsubsection{Oscillatori}
Vediamo come si può implementare un oscillatore per generare un segnale di clock. 

Il circuito che vediamo, nello specifico, è l'\textit{oscillatore ad anello}.
Questo si basa sull'utilizzo di 3 inverter (porte NOT) in serie.

\begin{center}
\begin{tikzpicture}[transform shape]
	% Paths, nodes and wires:
	\node[ieeestd not port](N1) at (4.985, 6){} node[anchor=center] at (N1.text){$1$};
	\node[ieeestd not port](N2) at (7, 6){} node[anchor=center] at (N2.text){$2$};
	\node[ieeestd not port](N3) at (8.985, 6){} node[anchor=center] at (N3.text){$3$};
	\draw (5.862, 6) -- (6.123, 6);
	\draw (7.877, 6) -- (8.108, 6);
	\draw (9.862, 6) -- (10, 6) -- (10, 4) -- (4, 4) -- (4, 6) -- (4.108, 6);
	\node[shape=rectangle, minimum width=0.465cm, minimum height=0.465cm] at (4, 6.5){} node[anchor=center, align=center, text width=0.077cm, inner sep=6pt] at (4, 6.5){$v_1$};
	\node[shape=rectangle, minimum width=0.465cm, minimum height=0.465cm] at (6, 6.5){} node[anchor=center, align=center, text width=0.077cm, inner sep=6pt] at (6, 6.5){$v_2$};
	\node[shape=rectangle, minimum width=0.465cm, minimum height=0.465cm] at (8, 6.5){} node[anchor=center, align=center, text width=0.077cm, inner sep=6pt] at (8, 6.5){$v_3$};
\end{tikzpicture}
\end{center}

Dobbiamo innanzitutto dimostrare che questo circuito è effettivamente astabile. Non c'è uno stato stabile da cui partire, per cui si sceglie uno stato qualsiasi: prendiamo quello dove $v_1$ è alta a $V_{DD}$. Poniamo quindi che il tempo di propagazione degli inverter sia pari a $\tau$. 

\begin{enumerate}
	\item Dopo la commutazione di $v_1$ a $V_{DD}$, il secondo inverter impiegherà un tempo $\tau$ ad "accorgersi" che $v_1$ è alto, quindi $v_2$ scenderà a $0 \, \mathrm{V}$;
	\item Dopo la commutazione di $v_2$ a $0$, il terzo inverter impiegherà un tempo $\tau$ ad "accorgersi" che $v_2$ è basso, quindi $v_3$ salirà a $V_{DD}$;
	\item Dopo la commutazione di $v_3$ a $V_{DD}$, il primo inverter impiegherà un tempo $\tau$ ad "accorgersi" che $v_3$ è alto, quindi $v_1$ scenderà a $0 \, \mathrm{V}$;
	\item Dopo la commutazione di $v_1$ a $0 \, \mathrm{V}$, il secondo inverter impiegherà un tempo $\tau$ ad "accorgersi" che $v_1$ è basso, quindi $v_2$ salirà a $V_{DD}$;
	\item Dopo la commutazione di $v_2$ a $V_{DD}$, il terzo inverter impiegherà un tempo $\tau$ ad "accorgersi" che $v_2$ è alto, quindi $v_3$ scenderà a $0 \, \mathrm{V}$;
	\item Dopo la commutazione di $v_3$ a $0$, il primo inverter impiegherà un tempo $\tau$ ad "accorgersi" che $v_3$ è basso, quindi $v_1$ salirà a $V_{DD}$;
\end{enumerate}

Dopo 6 iterazioni siamo quindi tornati alla situazione originale, e questo procedimento si ripeterà all'infinito, con periodo $6\tau$. Prendendo un numero qualsiasi di inverter (che deve essere dispari) $N$, abbiamo che il periodo è:
$$
T = 2 N \tau, \quad f = \frac{1}{2 N \tau}
$$
Il duty cycle sarà invece dato dal tempo che uno qualsiasi degli inverter passa alto sul periodo totale (la metà), cioè: 
$$
D = \frac{N \tau}{2 N \tau} = \frac{1}{2}
$$

Solitamente, poi, gli inverter (non reazionati) vengono usati come elementi di ritardo nei circuiti digitali, in quanto come abbiamo già visto in cascata introducono un ritardo uguale alla somma dei loro singoli ritardi.

\subsubsection{NE555}
Approfondiamo il funzionamento dei generatori di clock studiando un celebre circuito integrato con la funzione di timer, cioè il NE555.
Vediamone il funzionamento interno:
\begin{center}
\begin{tikzpicture}[transform shape]
	% Paths, nodes and wires:
	\draw (5, 8) to[american resistor, l={$R$}] (5, 6);
	\draw (5, 6) to[american resistor, l={$R$}] (5, 4);
	\draw (5, 4) to[american resistor, l={$R$}] (5, 2);
	\node[vcc] at (5, 8){};
	\node[ground] at (5, 1){};
	\node[op amp, yscale=-1] at (7.19, 6.49){};
	\node[op amp, yscale=-1] at (7.19, 3.51){};
	\draw (5, 4) -- (6, 4);
	\draw (5, 6) -- (6, 6);
	\draw (6, 6.98) -| (6, 8);
	\node[shape=rectangle, minimum width=0.465cm, minimum height=0.465cm] at (6, 8.5){} node[anchor=center, align=center, text width=0.077cm, inner sep=6pt] at (6, 8.5){$V_{TH}$};
	\draw (6, 3.02) -| (6, 2);
	\node[shape=rectangle, minimum width=0.465cm, minimum height=0.465cm] at (6, 1.5){} node[anchor=center, align=center, text width=0.077cm, inner sep=6pt] at (6, 1.5){$V_{TRIG}$};
	\node[flipflop SR, add async SR] at (10.13, 4.91){};
	\draw (8.38, 3.51) |- (8.5, 5.75);
	\draw (9.01, 4.07) -| (8.75, 6.25) |- (8.38, 6.49);
	\node[jump crossing] at (8.75, 5.75){};
	\draw (9, 5.75) -- (9, 5.75);
	\draw (8.61, 5.75) -- (8.5, 5.75);
	\draw (8.89, 5.75) -| (9.01, 5.75);
	\draw (10.13, 3.37) |- (10, 3);
	\node[shape=rectangle, minimum width=0.465cm, minimum height=0.465cm] at (9.5, 3){} node[anchor=center, align=center, text width=0.077cm, inner sep=6pt] at (9.5, 3){$\overline{R}$};
	\draw (12, 4) |- (11.25, 4.07);
	\draw (12, 4.07) -| (12, 5);
	\draw (12, 4) -| (12, 3);
	\draw (12, 3) to[american resistor, l={$R_1$}] (14, 3);
	\node[ieeestd not port] at (13, 5){};
	\draw (12, 5) -- (12.123, 5);
	\draw (13.877, 5) -- (14, 5);
	\node[npn] at (15, 3){};
	\draw (14.16, 3) |- (14, 3);
	\draw (14, 5) -- (15, 5);
	\draw (15, 3.77) -| (15, 4) -- (15.5, 4);
	\draw (15, 5) -- (15.5, 5);
	\node[shape=rectangle, minimum width=0.465cm, minimum height=0.465cm] at (16, 4){} node[anchor=center, align=center, text width=0.077cm, inner sep=6pt] at (16, 4){$D$};
	\node[shape=rectangle, minimum width=0.465cm, minimum height=0.465cm] at (16, 5){} node[anchor=center, align=center, text width=0.077cm, inner sep=6pt] at (16, 5){$Q$};
	\draw (5, 1) -| (5, 2);
	\draw (5, 1) -| (15, 2.23);
\end{tikzpicture}
\end{center}

Gli amplificatori differenziali in circuito aperto vengono usati come \textit{comparatori}, e l'integrato al centro è un latch SR. L'analisi del circuito ci porta a dire che:
\begin{itemize}
	\item Se $V_{TH} > \frac{2}{3} V_{CC}$, allora l'entrata di reset del latch ($R$) è posta a $1$;
	\item Se $V_{TR} < \frac{1}{3} V_{CC}$, allora l'entrata di set del latch ($S$) è posta a $1$.
\end{itemize}
La priorità è data all'$S$, per cui nello stato $R=1$, $S=1$, $Q$ è alta.

L'inverter in uscita a $Q$ è tarato per correnti maggiori di quelle gestite dal resto del circuito ($\sim 200 \, \mathrm{mA}$), per cui fa da stadio di uscita del circuito. Il $D$ viene detto \textit{Discharge}: se $\overline{Q} = 1$, il transistor è saturo e la $D$ viene posta fissa a $0 \, \mathrm{V}$. 

\par\smallskip
\noindent
\textbf{\textsf{Oscillatori col NE555}} \\

\noindent
La chiave nell'utilizzo dei circuiti integrati come il 555 sta nell'introduzione di catene di \textit{retroazione}. Dimostriamo tale fatto vedendo una delle configurazioni ad oscillatore tipiche:
\begin{center}
\begin{tikzpicture}[transform shape]
	% Paths, nodes and wires:
	\node[fourport](N1) at (8.09, 4.045){} node[anchor=south] at (N1.text){$\mathrm{NE555}$};
	\draw (9, 4.5) -- (10, 4.5);
	\draw (9, 3.59) -| (10.011, 3.592);
	\draw (7.18, 4.5) |- (6.75, 4.5) |- (7.18, 3.59);
	\draw (6, 4) to[american resistor, l={$R_A$}] (6, 6);
	\draw (6, 4) to[american resistor, l_={$R_B$}] (6, 2);
	\draw (6.75, 4) -- (6.5, 4) -- (6.5, 2) -- (6, 2);
	\node[vcc] at (6, 6){};
	\draw (6, 4) -| (4, 4) |- (10, -0) -| (10.011, 3.59);
	\draw (6.5, 2) to[capacitor, l={$C$}] (6.5, 1);
	\node[tlground] at (6.5, 1){};
	\node[tlground] at (6.497, 1.032){};
	\node[shape=rectangle, minimum width=0.465cm, minimum height=0.465cm] at (10.5, 4.5){} node[anchor=center, align=center, text width=0.077cm, inner sep=6pt] at (10.5, 4.5){$V_Q$};
	\node[shape=rectangle, minimum width=0.465cm, minimum height=0.465cm] at (7.65, 4.475){} node[anchor=center, align=center, text width=0.077cm, inner sep=6pt] at (7.35, 4.475){$TH$};
	\node[shape=rectangle, minimum width=0.465cm, minimum height=0.465cm] at (7.677, 3.632){} node[anchor=center, align=center, text width=0.077cm, inner sep=6pt] at (7.35, 3.632){$TR$};
	\node[shape=rectangle, minimum width=0.465cm, minimum height=0.465cm] at (8.58, 4.46){} node[anchor=center, align=center, text width=0.077cm, inner sep=6pt] at (8.58, 4.46){$Q$};
	\node[shape=rectangle, minimum width=0.465cm, minimum height=0.465cm] at (8.583, 3.63){} node[anchor=center, align=center, text width=0.077cm, inner sep=6pt] at (8.583, 3.63){$D$};
\end{tikzpicture}
\end{center}

Anche qui iniziamo a studiare il comportamento del circuito da uno stato qualsiasi, cioè quello ad uscita $Q$ alta.
\begin{enumerate}
\item \textbf{\textsf{Fase 1: set}} \\
	Se $Q = 1$ (cioè $V_{Q} = V_{CC}$) allora $D$ è in alta impedenza (in quanto significa che $\overline{Q}$ dell'SR è bassa). Possiamo quindi trascurare il feedback della linea di discharge. 

Il condensatore allora si caricherà, fino a raggiungere tensione fra le sue armature pari a $\frac{2}{3} V_{CC}$;

\item \textbf{\textsf{Fase 2: reset}} \\
La carica del condensatore, raggiunta $\frac{2}{3} V_{CC}$, porta l'entrata di reset del latch a 1. Visto che il set è alla stessa tensione, l'entrata di set è bassa, e quindi l'SR ha uscita bassa. Segue che la $\overline{Q}$ dell'SR è alta, per cui la linea di discharge è fissa a $0 \, \mathrm{V}$.

Il condensatore è quindi in scarica, finché non raggiunge la tensione $\frac{1}{3} V_{CC}$.

Una volta raggiunta la $\frac{1}{3} V_{CC}$, l'uscita del 555 commuta nuovamente, e ci troviamo nello stato iniziale. Questo permette di ottenere effettivamente un comportamento astabile.
\end{enumerate}

Calcoliamo un espressione per la $T_1$, cioè il tempo impiegato dalla capacità per passare da $\frac{1}{3} V_{CC}$ a $\frac{2}{3} V_{CC}$, e la $T_2$, cioè il tempo impiegato dalla capacità per passare nuovamente da $\frac{2}{3} V_{CC}$ a $\frac{1}{3} V_{CC}$. Abbiamo che la tensione sulla capacità è:
$$
V_C(t) = V_f + (V_i - V_f) e^{-\frac{t}{\tau}}
$$
nell'equazione si ha la tensione $V_{f}$ finale, da raggiungere ($0$ in scarica e $V_{CC}$ in carica), la tensione $V_i$ da cui si parte, e il tempo caratteristico del circuito RC $\tau$.

Quando è stata raggiunta $V_{com} = \frac{2}{3} V_{CC}$ (al tempo $T_1$), questa dà:
$$
V_C(t = T_1) = V_{com} = V_f + (V_i - V_f) e^{-\frac{T_1}{\tau}}
\, \Rightarrow \,
e^{\frac{T_1}{\tau}} = \frac{V_i - V_f}{V_{com} - V_f}
\, \Rightarrow \,
T_1 = \tau \ln \left( \frac{V_i - V_f}{V_{com} - V_f} \right)
$$

\begin{enumerate}
\item \textbf{\textsf{Fase 1: set}} \\
Nella prima fase, si ha che i parametri valgono:
$$
V_i = \frac{1}{3} V_{CC}, \quad V_f = V_{CC}, \quad V_{com} = \frac{2}{3} V_{CC}
$$
e la condizione da imporre è:
$$
V_i < V_{com} < V_f
$$

Resta da calcolare $\tau$, come:
$$
\tau = CR, \quad R = R_A + R_B \implies \tau = C (R_A + R_B)
$$
da cui si ha la $T_1$:
$$
T_1 = C (R_A + R_B) \ln \left( \frac{ \frac{1}{3} V_{CC} - V_{CC} }{ \frac{2}{3} V_{CC} - V_{CC} } \right) = C (R_A + R_B) \ln (2)
$$

\item \textbf{\textsf{Fase 2: reset}} \\
Nella seconda fase, si ha che i parametri valgono:
$$
V_i = \frac{2}{3} V_{CC}, \quad V_f = 0, \quad V_{com} = \frac{1}{3} V_{CC}
$$\begin{tikzpicture}[transform shape]
	% Paths, nodes and wires:
	\node[fourport](N1) at (8.09, 4.045){} node[anchor=south] at (N1.text){$\mathrm{NE555}$};
	\draw (9, 4.5) -- (10, 4.5);
	\draw (9, 3.59) -| (10.011, 3.592);
	\draw (7.18, 4.5) |- (6.75, 4.5) |- (7.18, 3.59);
	\draw (6, 4) to[american resistor, l={$R_A$}] (6, 6);
	\draw (6, 4) to[american resistor, l_={$R_B$}] (6, 2);
	\draw (6.75, 4) -- (6.5, 4) -- (6.5, 2) -- (6, 2);
	\node[vcc] at (6, 6){};
	\draw (6, 4) -| (4, 4) |- (10, -0) -| (10.011, 3.59);
	\draw (6.5, 2) to[capacitor, l={$C$}] (6.5, 1);
	\node[tlground] at (6.5, 1){};
	\node[tlground] at (6.497, 1.032){};
	\node[shape=rectangle, minimum width=0.465cm, minimum height=0.465cm] at (10.5, 4.5){} node[anchor=center, align=center, text width=0.077cm, inner sep=6pt] at (10.5, 4.5){$V_Q$};
	\node[shape=rectangle, minimum width=0.465cm, minimum height=0.465cm] at (7.65, 4.475){} node[anchor=center, align=center, text width=0.077cm, inner sep=6pt] at (7.65, 4.475){$TH$};
	\node[shape=rectangle, minimum width=0.465cm, minimum height=0.465cm] at (7.677, 3.632){} node[anchor=center, align=center, text width=0.077cm, inner sep=6pt] at (7.677, 3.632){$TR$};
	\node[shape=rectangle, minimum width=0.465cm, minimum height=0.465cm] at (8.58, 4.46){} node[anchor=center, align=center, text width=0.077cm, inner sep=6pt] at (8.58, 4.46){$Q$};
	\node[shape=rectangle, minimum width=0.465cm, minimum height=0.465cm] at (8.583, 3.63){} node[anchor=center, align=center, text width=0.077cm, inner sep=6pt] at (8.583, 3.63){$D$};
\end{tikzpicture}
e la condizione da imporre è:
$$
V_i > V_{com} > V_f
$$

La $\tau$ non è la stessa, in quanto una delle resistenze è a ground, e quindi:
$$
\tau = C R, \quad R = R_B, \quad \tau = C R_B
$$
da cui si ha la $T_2$:
$$
T_1 = C R_B \ln \left( \frac{ \frac{2}{3} V_{CC} - 0 }{ \frac{1}{3} V_{CC} - 0 } \right) = C R_B \ln (2)
$$
\end{enumerate}

Chiaramente abbiamo adesso un indicazione per il periodo complessivo $T$:
$$
T = T_1 + T_2 = C (R_A + 2 R_B) \ln (2)
$$
Ancora più interessanti sono le considerazioni che possiamo fare sul \textit{duty cycle}: questo varrà infatti:
$$
D = \frac{T_1}{T} = \frac{R_A + R_B}{R_A + 2 R_B}
$$
cioè il duty cycle non sarà mai veramente al $50 \%$, ma possiamo avvicinarci abbassando $R_B$.

\end{document}
